%!TEX TS-program = pdflatex

\documentclass[aspectratio=169,10pt]{beamer}

% -----------------------------------------------------------------------------
% Быстрая сборка: pdfLaTeX. Вся графика создаётся отдельным Python-скриптом.
% В Overleaf выберите Compiler = pdfLaTeX.
% -----------------------------------------------------------------------------
\usepackage[T2A]{fontenc}
\usepackage[utf8]{inputenc}
\usepackage[english,russian]{babel}
\usepackage{cmap}
\usepackage{amsmath,amssymb,mathtools}
\usepackage{booktabs}
\usepackage{array}
\usepackage{etoolbox}
\usepackage{PTSans}
\usepackage{graphicx}
\graphicspath{{figures/}}

\renewcommand{\familydefault}{\sfdefault}
\renewcommand{\rmdefault}{\sfdefault}

\newbool{russian}
\booltrue{russian}
\usepackage{HSE-theme/beamerthemeHSE}

\hypersetup{bookmarks=false}
\setbeamertemplate{blocks}[rounded][shadow=false]
\setbeamerfont{frametitle}{size=\Large}
\setbeamerfont{framesubtitle}{size=\small}
\setbeamerfont{block title}{size=\normalsize}
\setbeamerfont{block body}{size=\small}

\newcommand{\R}{\mathbb{R}}
\newcommand{\norm}[1]{\left\lVert #1\right\rVert}
\newcommand{\ip}[2]{\left\langle #1,#2\right\rangle}
\newcommand{\grad}{\nabla}
\newcommand{\Hess}{\nabla^2}
\newcommand{\argmin}{\operatorname*{arg\,min}}
\newcommand{\lammin}{\lambda_{\min}}
\newcommand{\lammax}{\lambda_{\max}}

\title[Методы оптимизации в ML]{Методы оптимизации\\в машинном обучении}
\subtitle{Лекция 3. Квадратичная оптимизация: геометрия, спектр и точное решение}
\author[Н. Лукьяненко]{Никита Лукьяненко}
\institute{Совместная программа СМОЛ ГУ и ФКН НИУ ВШЭ}
\date{2026/27 учебный год}

\begin{document}

% 1 ---------------------------------------------------------------------------
\frame[plain]{\titlepage}

% 2 ---------------------------------------------------------------------------
\begin{frame}{От локальной геометрии --- к точному классу задач}
\small
На прошлой лекции мы получили локальную квадратичную модель
\[
  f(x+h)\approx f(x)+\grad f(x)^\top h+\frac12 h^\top \Hess f(x)h.
\]

\begin{columns}[T,onlytextwidth]
\column{0.46\textwidth}
\begin{block}{Первый порядок}
\[
\grad f(x)
\]
описывает локальный наклон.
\end{block}

\column{0.46\textwidth}
\begin{block}{Второй порядок}
\[
\Hess f(x)
\]
описывает локальную кривизну.
\end{block}
\end{columns}

\vspace{0.3em}
Сегодня рассмотрим класс функций, для которых квадратичная модель описывает функцию \textbf{точно}.
\end{frame}

% 3 ---------------------------------------------------------------------------
\begin{frame}{Когда квадратичное приближение становится точным?}
\small
\begin{columns}[T,onlytextwidth]
\column{0.43\textwidth}
Для общей гладкой функции формула второго порядка лишь приближает поведение около выбранной точки.

\vspace{0.7em}
Для квадратичной функции остатка нет:
\[
  f(x+h)=f(x)+\grad f(x)^\top h+\frac12 h^\top \Hess f(x)h.
\]

\begin{alertblock}{Почему это удобно?}
Мы можем полностью понять геометрию задачи, а затем использовать её как модель для более общих функций.
\end{alertblock}

\column{0.53\textwidth}
\centering
\includegraphics[width=0.98\linewidth]{quadratic_taylor_exact.pdf}
\end{columns}
\end{frame}

% 4 ---------------------------------------------------------------------------
\begin{frame}{Почему квадратичные задачи важны?}
\vspace{0.45em}

\begin{block}{1. Их можно исследовать полностью}
Условия существования и единственности решения выражаются через спектр одной матрицы.
\end{block}

\begin{block}{2. Они описывают локальную геометрию}
Около точки минимума гладкая функция во втором порядке выглядит квадратично.
\end{block}

\begin{block}{3. Они возникают непосредственно в машинном обучении}
Метод наименьших квадратов --- классический пример квадратичной оптимизации.
\end{block}
\end{frame}

% 5 ---------------------------------------------------------------------------
\begin{frame}{Одна формула --- очень разная геометрия}
\small
\centering
\includegraphics[width=0.92\linewidth]{three_quadratic_geometries.pdf}

\vspace{0.15em}
\begin{alertblock}{Главный вопрос лекции}
\centering
Как по матрице квадратичной функции заранее понять, какая геометрия перед нами?
\end{alertblock}
\end{frame}

% 6 ---------------------------------------------------------------------------
\begin{frame}{Квадратичная функция}
\vspace{0.45em}
\begin{block}{Определение}
Функция
\[
  \boxed{
  f(x)=\frac12 x^\top A x-b^\top x+c,
  \qquad x\in\R^d
  }
\]
называется квадратичной, где
\[
  A\in\R^{d\times d},\qquad b\in\R^d,\qquad c\in\R.
\]
\end{block}

\vspace{0.35em}
Матрица $A$ будет отвечать за кривизну и форму поверхности; $b$ --- за положение её центра; $c$ --- за вертикальный сдвиг.
\end{frame}

% 7 ---------------------------------------------------------------------------
\begin{frame}{Что делает каждый член?}
\small
\begin{columns}[T,onlytextwidth]
\column{0.36\textwidth}
\[
 f(x)=\frac12x^\top Ax-b^\top x+c.
\]

\begin{itemize}
  \item $A$ задаёт \textbf{форму};
  \item $b$ сдвигает оптимум;
  \item $c$ не меняет оптимум вообще.
\end{itemize}

\begin{block}{Важно}
Изменяя $b$ и $c$, мы не меняем гессиан квадратичной функции.
\end{block}

\column{0.60\textwidth}
\centering
\includegraphics[width=0.99\linewidth]{quadratic_terms_effect.pdf}
\end{columns}
\end{frame}

% 8 ---------------------------------------------------------------------------
\begin{frame}{Зачем перед квадратичной формой ставят $1/2$?}
\vspace{0.65em}
Множитель $1/2$ не меняет сути задачи. Он выбран ради удобства производных.

\vspace{0.55em}
В одном измерении:
\[
  \frac{d}{dx}\left(\frac12 ax^2\right)=ax.
\]

В нескольких измерениях при симметричной $A$ аналогично:
\[
  \grad\left(\frac12x^\top Ax\right)=Ax.
\]

\begin{alertblock}{Соглашение}
В оптимизации запись $\frac12x^\top Ax$ используется очень часто именно потому, что она убирает лишнюю двойку из градиента.
\end{alertblock}
\end{frame}

% 9 ---------------------------------------------------------------------------
\begin{frame}{Обязательно ли считать $A$ симметричной?}
\small
Любую квадратную матрицу можно разложить как
\[
 A=\underbrace{\frac{A+A^\top}{2}}_{S}
 +\underbrace{\frac{A-A^\top}{2}}_{K},
\]
где
\[
 S^\top=S,\qquad K^\top=-K.
\]

\begin{block}{Вопрос}
Может ли кососимметричная часть $K$ влиять на величину $x^\top Ax$?
\end{block}

На следующем слайде ответим на этот вопрос формально.
\end{frame}

% 10 --------------------------------------------------------------------------
\begin{frame}{Кососимметричная часть не влияет на квадратичную форму}
\small
\begin{block}{Утверждение}
Если $K^\top=-K$, то для любого $x\in\R^d$
\[
  \boxed{x^\top Kx=0.}
\]
\end{block}

\begin{block}{Доказательство}
Величина $x^\top Kx$ --- число, поэтому она совпадает со своим транспонированием:
\[
 x^\top Kx=(x^\top Kx)^\top=x^\top K^\top x.
\]
Но $K^\top=-K$, следовательно,
\[
 x^\top Kx=-x^\top Kx.
\]
Значит,
\[
 2x^\top Kx=0,
\]
что и требовалось доказать.
\end{block}
\end{frame}

% 11 --------------------------------------------------------------------------
\begin{frame}{Поэтому можно считать $A$ симметричной}
\small
Из разложения $A=S+K$ и предыдущего утверждения получаем
\[
  x^\top Ax=x^\top Sx+x^\top Kx=x^\top Sx.
\]
Следовательно,
\[
  \boxed{
  x^\top Ax=x^\top\frac{A+A^\top}{2}x.
  }
\]

\begin{alertblock}{Вывод}
Для квадратичной оптимизации несимметричная часть матрицы не видна вообще. Поэтому дальше без ограничения общности считаем
\[
  \boxed{A=A^\top.}
\]
\end{alertblock}
\end{frame}

% 12 --------------------------------------------------------------------------
\begin{frame}{Пример в $\R^2$}
\small
\begin{columns}[T,onlytextwidth]
\column{0.45\textwidth}
Пусть
\[
A=\begin{pmatrix}2&1\\1&4\end{pmatrix},
\qquad
b=\begin{pmatrix}1\\2\end{pmatrix}.
\]
Тогда при $c=0$
\[
\begin{aligned}
f(x_1,x_2)
&=\frac12x^\top Ax-b^\top x\\
&=x_1^2+x_1x_2+2x_2^2-x_1-2x_2.
\end{aligned}
\]

Линии уровня уже дают важную информацию о форме задачи.

\column{0.51\textwidth}
\centering
\includegraphics[width=0.92\linewidth]{quadratic_example_contours.pdf}
\end{columns}
\end{frame}

% 13 --------------------------------------------------------------------------
\begin{frame}{Самопроверка: восстановите матричную запись}
\small
Рассмотрим
\[
 f(x,y)=3x^2-2xy+y^2+4x-6y+7.
\]

Представьте её в форме
\[
 f(z)=\frac12z^\top Az-b^\top z+c,
 \qquad
 z=\begin{pmatrix}x\\y\end{pmatrix}.
\]

\vspace{0.5em}
\begin{block}{Вопросы}
Каковы $A$, $b$ и $c$? Почему коэффициент перед $xy$ нужно особенно внимательно сопоставлять с матрицей?
\end{block}

\vfill
\centering
\textbf{На обсуждение: около минуты.}
\end{frame}

% 14 --------------------------------------------------------------------------
\begin{frame}{Градиент квадратичной формы}
\small
\begin{block}{Утверждение}
Пусть $A=A^\top$. Тогда
\[
 \boxed{
 \grad\left(\frac12x^\top Ax\right)=Ax.
 }
\]
\end{block}

Это многомерный аналог знакомой формулы
\[
 \frac{d}{dx}\left(\frac12ax^2\right)=ax.
\]

\begin{alertblock}{Почему доказательство полезно?}
Оно показывает, как аккуратно работать с индексами и где именно используется симметрия $A$.
\end{alertblock}
\end{frame}

% 15 --------------------------------------------------------------------------
\begin{frame}{Доказательство формулы для градиента}
\footnotesize
Запишем
\[
 \frac12x^\top Ax=\frac12\sum_{i=1}^{d}\sum_{j=1}^{d}A_{ij}x_ix_j.
\]
Для $k$-й координаты градиента
\[
\begin{aligned}
\frac{\partial}{\partial x_k}\left(\frac12x^\top Ax\right)
&=\frac12\sum_{j=1}^{d}A_{kj}x_j
 +\frac12\sum_{i=1}^{d}A_{ik}x_i.
\end{aligned}
\]
Так как $A_{ik}=A_{ki}$,
\[
\begin{aligned}
\frac{\partial}{\partial x_k}\left(\frac12x^\top Ax\right)
&=\sum_{j=1}^{d}A_{kj}x_j
=(Ax)_k.
\end{aligned}
\]
Это верно для всех $k$, поэтому
\[
 \boxed{\grad\left(\frac12x^\top Ax\right)=Ax.}
\]
\end{frame}

% 16 --------------------------------------------------------------------------
\begin{frame}{Линейный и постоянный члены}
\vspace{0.6em}
Для линейного члена
\[
 b^\top x=\sum_{i=1}^{d}b_ix_i,
\]
поэтому
\[
 \grad(b^\top x)=b.
\]
Следовательно,
\[
 \grad(-b^\top x)=-b.
\]
А постоянный член не зависит от $x$:
\[
 \grad c=0.
\]

\begin{block}{Итого}
Для квадратичной функции остаётся только сложить три простых результата.
\end{block}
\end{frame}

% 17 --------------------------------------------------------------------------
\begin{frame}{Главная формула: градиент квадратичной функции}
\vspace{0.7em}
Для
\[
 f(x)=\frac12x^\top Ax-b^\top x+c,
 \qquad A=A^\top,
\]
получаем
\[
 \boxed{\grad f(x)=Ax-b.}
\]

\begin{block}{Обратите внимание}
Градиент линейно зависит от текущей точки $x$. Это сделает анализ градиентного спуска на следующей лекции почти полностью алгебраическим.
\end{block}
\end{frame}

% 18 --------------------------------------------------------------------------
\begin{frame}{Гессиан ещё проще}
\vspace{0.7em}
Поскольку
\[
 \grad f(x)=Ax-b,
\]
ещё одно дифференцирование даёт
\[
 \boxed{\Hess f(x)=A.}
\]

\begin{alertblock}{Ключевая особенность}
Гессиан квадратичной функции не зависит от точки $x$. Кривизна поверхности одинакова везде.
\end{alertblock}
\end{frame}

% 19 --------------------------------------------------------------------------
\begin{frame}{У общей функции кривизна меняется, у квадратичной --- нет}
\small
\begin{columns}[T,onlytextwidth]
\column{0.38\textwidth}
В одном измерении кривизну отражает вторая производная.

\vspace{0.4em}
Для общей функции:
\[
 f''(x)\ \text{зависит от }x.
\]

Для квадратичной:
\[
 q''(x)=\text{const}.
\]

В многомерном случае роль второй производной играет матрица $\Hess f(x)$.

\column{0.58\textwidth}
\centering
\includegraphics[width=0.98\linewidth]{constant_curvature_quadratic.pdf}
\end{columns}
\end{frame}

% 20 --------------------------------------------------------------------------
\begin{frame}{Кривизна вдоль направления}
\small
С прошлой лекции мы знаем: для единичного направления $v$
\[
 \frac{d^2}{dt^2}f(x+tv)\bigg|_{t=0}
 =v^\top\Hess f(x)v.
\]
Для квадратичной функции
\[
 \boxed{
 \kappa(v)=v^\top Av.
 }
\]

\begin{block}{Что особенно удобно}
Эта величина зависит от направления $v$, но \textbf{не зависит от точки} $x$.
\end{block}

Теперь вопрос о геометрии функции превращается в вопрос о квадратичной форме $v^\top Av$.
\end{frame}

% 21 --------------------------------------------------------------------------
\begin{frame}{Начнём с диагональной матрицы}
\small
Пусть
\[
 A=\begin{pmatrix}
 \lambda_1&0\\
 0&\lambda_2
 \end{pmatrix}.
\]
Тогда
\[
 f(x)=\frac12x^\top Ax
 =\frac12\left(\lambda_1x_1^2+\lambda_2x_2^2\right).
\]

\begin{block}{Преимущество диагонального случая}
Координаты не взаимодействуют: вклад каждого направления виден отдельно.
\end{block}

Именно к такой форме мы сведём произвольную симметричную матрицу с помощью собственных векторов.
\end{frame}

% 22 --------------------------------------------------------------------------
\begin{frame}{Если собственные значения одинаковы}
\small
\begin{columns}[T,onlytextwidth]
\column{0.43\textwidth}
При
\[
 A=I
\]
имеем
\[
 f(x)=\frac12(x_1^2+x_2^2).
\]
Линия уровня $f(x)=C$ задаётся условием
\[
 x_1^2+x_2^2=2C,
\]
то есть является окружностью.

\begin{block}{Геометрия}
Кривизна одинакова во всех направлениях.
\end{block}

\column{0.53\textwidth}
\centering
\includegraphics[width=0.88\linewidth]{contours_equal_eigs.pdf}
\end{columns}
\end{frame}

% 23 --------------------------------------------------------------------------
\begin{frame}{Если собственные значения различаются}
\small
\begin{columns}[T,onlytextwidth]
\column{0.43\textwidth}
Для
\[
 A=\begin{pmatrix}1&0\\0&10\end{pmatrix}
\]
получаем
\[
 f(x)=\frac12(x_1^2+10x_2^2).
\]
Линии уровня становятся эллипсами.

\begin{block}{Интуиция}
В направлении $x_2$ функция растёт значительно быстрее, поэтому допустимое отклонение по этой координате меньше.
\end{block}

\column{0.53\textwidth}
\centering
\includegraphics[width=0.88\linewidth]{contours_unequal_eigs.pdf}
\end{columns}
\end{frame}

% 24 --------------------------------------------------------------------------
\begin{frame}{Что означает большое собственное значение?}
\small
Если $v_i$ --- собственный вектор с собственным значением $\lambda_i$, то
\[
 Av_i=\lambda_i v_i.
\]
Рассмотрим движение только вдоль этого направления: $x=tv_i$. Для $\norm{v_i}=1$
\[
 f(tv_i)=\frac12t^2v_i^\top Av_i
 =\frac12\lambda_i t^2.
\]

\begin{columns}[T,onlytextwidth]
\column{0.45\textwidth}
\begin{alertblock}{Вывод}
Чем больше $\lambda_i$, тем сильнее функция изгибается и быстрее растёт вдоль направления $v_i$.
\end{alertblock}
\column{0.51\textwidth}
\centering
\includegraphics[width=0.92\linewidth]{eigenvalue_curvature_sections.pdf}
\end{columns}
\end{frame}

% 25 --------------------------------------------------------------------------
\begin{frame}{Недиагональная матрица просто поворачивает главные оси}
\small
\begin{columns}[T,onlytextwidth]
\column{0.42\textwidth}
Например,
\[
 A=\begin{pmatrix}5&4\\4&5\end{pmatrix}.
\]
Линии уровня по-прежнему эллипсы, но их оси уже не совпадают с координатными осями.

\begin{block}{Новый вопрос}
Как найти естественные направления, вдоль которых квадратичная функция снова выглядит как сумма независимых квадратов?
\end{block}

Ответ --- собственные векторы $A$.

\column{0.54\textwidth}
\centering
\includegraphics[width=0.94\linewidth]{rotated_contours_eigenvectors.pdf}
\end{columns}
\end{frame}

% 26 --------------------------------------------------------------------------
\begin{frame}{Спектральная теорема для симметричной матрицы}
\small
\begin{block}{Теорема}
Если $A=A^\top$, то существует ортогональная матрица $Q$ и диагональная матрица $\Lambda$, такие что
\[
 \boxed{A=Q\Lambda Q^\top,}
\]
где столбцы $Q$ --- ортонормированные собственные векторы $A$, а
\[
 \Lambda=\operatorname{diag}(\lambda_1,\ldots,\lambda_d).
\]
\end{block}

\begin{alertblock}{На этой лекции}
Спектральную теорему используем без доказательства. Нас интересуют её последствия для оптимизации.
\end{alertblock}
\end{frame}

% 27 --------------------------------------------------------------------------
\begin{frame}{Переходим в собственные координаты}
\small
Положим
\[
 z=Q^\top x,
 \qquad\text{то есть}\qquad
 x=Qz.
\]
Тогда
\[
 x^\top Ax
 =z^\top Q^\top(Q\Lambda Q^\top)Qz
 =z^\top\Lambda z.
\]
Следовательно,
\[
 \boxed{
 x^\top Ax=\sum_{i=1}^{d}\lambda_i z_i^2.
 }
\]

\begin{block}{Смысл}
После поворота координат любая симметричная квадратичная форма становится суммой независимых одномерных квадратичных функций.
\end{block}
\end{frame}

% 28 --------------------------------------------------------------------------
\begin{frame}{Почему преобразование действительно является только поворотом?}
\small
Матрица $Q$ ортогональна:
\[
 Q^\top Q=I.
\]
Поэтому она сохраняет длины:
\[
 \norm{Qz}^2=z^\top Q^\top Qz=\norm{z}^2.
\]
И сохраняет скалярные произведения:
\[
 (Qz_1)^\top(Qz_2)=z_1^\top z_2.
\]

\begin{alertblock}{Геометрическая интерпретация}
Переход $x=Qz$ не растягивает пространство. Он лишь меняет ортонормированный базис так, чтобы его оси совпали с собственными направлениями $A$.
\end{alertblock}
\end{frame}

% 29 --------------------------------------------------------------------------
\begin{frame}{Как из собственных значений увидеть форму линии уровня?}
\footnotesize
В собственных координатах возьмём
\[
 f(z)=\frac12\left(z_1^2+6z_2^2\right),
 \qquad
 \lambda_1=1,\quad \lambda_2=6.
\]
На линии уровня $f(z)=1$
\[
 z_1^2+6z_2^2=2,
 \qquad
 a_1=\sqrt{\frac{2}{\lambda_1}}=\sqrt2,
 \quad
 a_2=\sqrt{\frac{2}{\lambda_2}}=\sqrt{\frac13}.
\]

\centering
\includegraphics[width=0.72\linewidth]{spectrum_geometry.pdf}

\vspace{-0.2em}
\textbf{Вывод:} большее $\lambda_i$ означает более сильную кривизну вдоль $v_i$ и более короткую полуось линии уровня.
\end{frame}

% 30 --------------------------------------------------------------------------
\begin{frame}{Спектральные границы квадратичной формы}
\small
\begin{block}{Утверждение}
Для симметричной $A$ и любого $x\in\R^d$
\[
 \boxed{
 \lammin(A)\norm{x}^2
 \le x^\top Ax
 \le \lammax(A)\norm{x}^2.
 }
\]
\end{block}

\begin{block}{Доказательство}
В собственных координатах $z=Q^\top x$:
\[
 x^\top Ax=\sum_i\lambda_i z_i^2.
\]
Так как $\lammin\le\lambda_i\le\lammax$,
\[
 \lammin\sum_i z_i^2
 \le\sum_i\lambda_i z_i^2
 \le\lammax\sum_i z_i^2.
\]
Ортогональность $Q$ даёт $\norm{z}=\norm{x}$.
\end{block}
\end{frame}

% 31 --------------------------------------------------------------------------
\begin{frame}{Положительная определённость}
\small
\begin{block}{Определение}
Симметричная матрица $A$ называется \textbf{положительно определённой}, если
\[
 \boxed{x^\top Ax>0\quad\text{для любого }x\ne0.}
\]
Обозначение:
\[
 A\succ0.
\]
\end{block}

Если допускается равенство нулю,
\[
 x^\top Ax\ge0\quad\forall x,
\]
то $A$ называется положительно полуопределённой и пишут
\[
 A\succeq0.
\]
\end{frame}

% 32 --------------------------------------------------------------------------
\begin{frame}{Как проверить $A\succ0$, не перебирая все направления?}
\small
Определение требует проверять знак $x^\top Ax$ для бесконечного множества векторов $x$. Для симметричной матрицы спектр превращает эту проверку в конечную.

\begin{block}{Спектральный критерий}
Для симметричной матрицы $A$
\[
 \boxed{
 A\succ0
 \quad\Longleftrightarrow\quad
 \lambda_{\min}(A)>0.
 }
\]
То есть достаточно посмотреть на \textbf{наименьшее} собственное значение.
\end{block}

\vspace{0.25em}
\begin{itemize}
  \item $\lambda_{\min}>0$ --- положительно определённая чаша;
  \item $\lambda_{\min}=0$ и отрицательных значений нет --- есть плоские направления;
  \item $\lambda_{\min}<0$ --- существует направление отрицательной кривизны.
\end{itemize}

Следующий слайд --- доказательство критерия.
\end{frame}

% 33 --------------------------------------------------------------------------
\begin{frame}{Доказательство спектрального критерия}
\footnotesize
Пусть сначала все $\lambda_i>0$. В собственных координатах
\[
 x^\top Ax=\sum_i\lambda_i z_i^2.
\]
Если $x\ne0$, то и $z\ne0$, поэтому хотя бы одно $z_i\ne0$, а значит
\[
 x^\top Ax>0.
\]
Следовательно, $A\succ0$.

\vspace{0.5em}
Теперь пусть $A\succ0$. Возьмём собственный вектор $v_i\ne0$:
\[
 Av_i=\lambda_i v_i.
\]
Тогда
\[
 v_i^\top Av_i
 =\lambda_i\norm{v_i}^2.
\]
Левая часть положительна, а $\norm{v_i}^2>0$. Поэтому $\lambda_i>0$.
\[
 \boxed{A\succ0\iff \lambda_i>0\ \forall i.}
\]
\end{frame}

% 34 --------------------------------------------------------------------------
\begin{frame}{$A\succ0$: единственная чашеобразная геометрия}
\small
\begin{columns}[T,onlytextwidth]
\column{0.42\textwidth}
Для
\[
 f(x)=\frac12x^\top Ax,
 \qquad A\succ0,
\]
имеем
\[
 f(x)>0\quad\text{при }x\ne0,
\qquad
 f(0)=0.
\]
Поэтому
\[
 \boxed{x^\star=0}
\]
--- единственный глобальный минимум.

\begin{block}{Геометрия}
Во всех направлениях кривизна положительна: поверхность изгибается вверх.
\end{block}

\column{0.54\textwidth}
\centering
\includegraphics[width=0.90\linewidth]{pd_bowl.pdf}
\end{columns}
\end{frame}

% 35 --------------------------------------------------------------------------
\begin{frame}{Отрицательное собственное значение разрушает ограниченность снизу}
\small
Пусть у $A$ есть собственный вектор $v$ с
\[
 Av=\lambda v,
 \qquad \lambda<0.
\]
Рассмотрим точки $x=tv$. Тогда
\[
 f(tv)=\frac12t^2v^\top Av
 =\frac12\lambda t^2\norm{v}^2.
\]
Так как $\lambda<0$,
\[
 f(tv)\longrightarrow -\infty
 \qquad\text{при }|t|\to\infty.
\]

\begin{alertblock}{Вывод}
Если у квадратичной части есть отрицательное собственное значение, то чистая квадратичная функция $\frac12x^\top Ax$ не ограничена снизу.
\end{alertblock}
\end{frame}

% 36 --------------------------------------------------------------------------
\begin{frame}{Собственные значения разных знаков: седловая геометрия}
\small
\begin{columns}[T,onlytextwidth]
\column{0.42\textwidth}
Если существуют
\[
 \lambda_i>0,
 \qquad
 \lambda_j<0,
\]
то в одном собственном направлении функция растёт вверх, а в другом --- вниз.

\begin{block}{Следствие}
Точка $x=0$ стационарна для $\frac12x^\top Ax$, но не является ни минимумом, ни максимумом.
\end{block}

Такую точку называют седловой.

\column{0.54\textwidth}
\centering
\includegraphics[width=0.90\linewidth]{indefinite_saddle.pdf}
\end{columns}
\end{frame}

% 37 --------------------------------------------------------------------------
\begin{frame}{Полуопределённость: минимум может быть не единственным}
\small
\begin{columns}[T,onlytextwidth]
\column{0.42\textwidth}
Пусть
\[
 A=\begin{pmatrix}1&0\\0&0\end{pmatrix}.
\]
Тогда
\[
 f(x_1,x_2)=\frac12x_1^2.
\]
Минимальное значение равно нулю для всех точек
\[
 x_1=0.
\]

\begin{alertblock}{Важно}
Нулевое собственное значение означает направление с нулевой кривизной. В этом направлении поверхность плоская.
\end{alertblock}

\column{0.54\textwidth}
\centering
\includegraphics[width=0.90\linewidth]{psd_valley.pdf}
\end{columns}
\end{frame}

% 38 --------------------------------------------------------------------------
\begin{frame}{Но при $A\succeq0$ линейный член может всё изменить}
\small
Рассмотрим
\[
 f(x,y)=x^2-y.
\]
Её гессиан
\[
 \Hess f=\begin{pmatrix}2&0\\0&0\end{pmatrix}\succeq0.
\]
Но при $y\to+\infty$
\[
 f(x,y)\to-\infty.
\]

\begin{alertblock}{Вывод}
Положительная полуопределённость квадратичной части сама по себе ещё не гарантирует существование минимума, если присутствует линейный член.
\end{alertblock}

Общий полуопределённый случай пока оставим за рамками курса. Для основной теории сосредоточимся на $A\succ0$.
\end{frame}

% 39 --------------------------------------------------------------------------
\begin{frame}{Ищем минимум при $A\succ0$}
\small
Для
\[
 f(x)=\frac12x^\top Ax-b^\top x+c
\]
имеем
\[
 \grad f(x)=Ax-b.
\]
В точке внутреннего минимума необходимо
\[
 \grad f(x^\star)=0.
\]
Следовательно,
\[
 \boxed{Ax^\star=b.}
\]

То есть задача минимизации свелась к решению системы линейных уравнений.
\end{frame}

% 40 --------------------------------------------------------------------------
\begin{frame}{Точное решение}
\vspace{0.8em}
Если $A\succ0$, то $A$ обратима. Поэтому из
\[
 Ax^\star=b
\]
следует
\[
 \boxed{x^\star=A^{-1}b.}
\]

\begin{alertblock}{Но есть важный практический нюанс}
В вычислениях почти никогда не стоит явно находить $A^{-1}$. Обычно решают систему
\[
 Ax=b
\]
специализированным численным методом.
\end{alertblock}

Формула $A^{-1}b$ важна для теории, но не обязательно является хорошим алгоритмом.
\end{frame}

% 41 --------------------------------------------------------------------------
\begin{frame}{Почему положительно определённая матрица обратима?}
\small
\begin{block}{Утверждение}
Если $A\succ0$, то $A$ невырождена.
\end{block}

\begin{block}{Доказательство}
Предположим, что
\[
 Ax=0.
\]
Тогда
\[
 x^\top Ax=x^\top0=0.
\]
Но положительная определённость требует
\[
 x^\top Ax>0
\]
для любого $x\ne0$. Значит, равенство $Ax=0$ возможно только при $x=0$.

Следовательно,
\[
 \ker A=\{0\},
\]
а значит $A$ обратима.
\end{block}
\end{frame}

% 42 --------------------------------------------------------------------------
\begin{frame}{Теорема о минимуме квадратичной функции}
\small
\begin{block}{Теорема}
Пусть
\[
 f(x)=\frac12x^\top Ax-b^\top x+c,
 \qquad A=A^\top\succ0.
\]
Тогда $f$ имеет единственный глобальный минимум, причём
\[
 \boxed{x^\star=A^{-1}b.}
\]
\end{block}

\begin{block}{Почему это сильнее условия $\grad f=0$?}
Условие первого порядка само по себе говорит лишь о стационарности. Положительная определённость позволяет доказать и глобальность, и единственность минимума.
\end{block}
\end{frame}

% 43 --------------------------------------------------------------------------
\begin{frame}{Доказательство: дополняем до квадрата}
\footnotesize
Пусть $x^\star=A^{-1}b$, то есть $b=Ax^\star$. Рассмотрим разность
\[
 f(x)-f(x^\star).
\]
Подставляя $b=Ax^\star$ и используя симметрию $A$, получаем
\[
\begin{aligned}
f(x)-f(x^\star)
&=\frac12x^\top Ax-(x^\star)^\top Ax
  +\frac12(x^\star)^\top Ax^\star\\
&=\frac12(x-x^\star)^\top A(x-x^\star).
\end{aligned}
\]
Так как $A\succ0$,
\[
 (x-x^\star)^\top A(x-x^\star)>0
 \qquad\text{для }x\ne x^\star.
\]
Следовательно,
\[
 f(x)>f(x^\star)\qquad\forall x\ne x^\star.
\]
Значит, $x^\star$ --- единственный глобальный минимум.
\end{frame}

% 44 --------------------------------------------------------------------------
\begin{frame}{Линейный член сдвигает минимум, но не меняет форму}
\small
\begin{columns}[T,onlytextwidth]
\column{0.39\textwidth}
При фиксированной $A\succ0$
\[
 \Hess f(x)=A
\]
не зависит от $b$.

Но положение минимума
\[
 x^\star=A^{-1}b
\]
зависит от $b$.

\begin{alertblock}{Геометрически}
$b$ передвигает центр эллиптической чаши, не меняя её главные направления и кривизну.
\end{alertblock}

\column{0.57\textwidth}
\centering
\includegraphics[width=0.98\linewidth]{linear_term_shifts_center.pdf}
\end{columns}
\end{frame}

% 45 --------------------------------------------------------------------------
\begin{frame}{Связь с линейной регрессией}
\small
В курсе машинного обучения вы уже работаете с задачей наименьших квадратов. Здесь посмотрим только на её \textbf{оптимизационную структуру}.

\vspace{0.5em}
Пусть
\[
 X\in\R^{n\times d},\qquad y\in\R^n,
\]
а предсказания имеют вид $Xw$.

Среднеквадратичная задача:
\[
 \boxed{
 \min_w\frac12\norm{Xw-y}^2.
 }
\]

Это квадратичная оптимизация по параметрам $w$.
\end{frame}

% 46 --------------------------------------------------------------------------
\begin{frame}{Разворачиваем квадрат нормы}
\small
Используем
\[
 \norm{Xw-y}^2=(Xw-y)^\top(Xw-y).
\]
Тогда
\[
\begin{aligned}
f(w)
&=\frac12(Xw-y)^\top(Xw-y)\\
&=\frac12w^\top X^\top Xw-y^\top Xw+\frac12y^\top y.
\end{aligned}
\]

\begin{block}{Сравним с общей формой}
\[
 \frac12w^\top Aw-b^\top w+c.
\]
\end{block}
\end{frame}

% 47 --------------------------------------------------------------------------
\begin{frame}{Матрица квадратичной задачи для линейной регрессии}
\vspace{0.65em}
Из предыдущего разложения:
\[
 \boxed{A=X^\top X,}
 \qquad
 \boxed{b=X^\top y,}
 \qquad
 \boxed{c=\frac12y^\top y.}
\]
Следовательно,
\[
 \boxed{\grad f(w)=X^\top(Xw-y),}
\]
\[
 \boxed{\Hess f(w)=X^\top X.}
\]

\begin{alertblock}{Оптимизационная геометрия линейной регрессии}
Она полностью определяется спектром матрицы $X^\top X$.
\end{alertblock}
\end{frame}

% 48 --------------------------------------------------------------------------
\begin{frame}{Почему $X^\top X$ всегда положительно полуопределена?}
\small
\begin{block}{Утверждение}
Для любой матрицы $X$
\[
 X^\top X\succeq0.
\]
\end{block}

\begin{block}{Доказательство}
Для любого $z\in\R^d$
\[
 z^\top X^\top Xz=(Xz)^\top(Xz)=\norm{Xz}^2\ge0.
\]
Следовательно, по определению
\[
 \boxed{X^\top X\succeq0.}
\]
\end{block}

Это объясняет, почему функция наименьших квадратов никогда не имеет седловой квадратичной геометрии.
\end{frame}

% 49 --------------------------------------------------------------------------
\begin{frame}{Когда решение линейной регрессии единственно?}
\small
Для $X^\top X$ положительная определённость означает
\[
 z^\top X^\top Xz=\norm{Xz}^2>0
 \qquad\forall z\ne0.
\]
Это эквивалентно условию
\[
 Xz\ne0\qquad\forall z\ne0,
\]
то есть
\[
 \ker X=\{0\}.
\]

\begin{alertblock}{Иными словами}
Если столбцы $X$ линейно независимы, то $X^\top X\succ0$ и решение единственно.
\end{alertblock}

Тогда условие $\grad f(w)=0$ даёт нормальные уравнения
\[
 X^\top Xw=X^\top y.
\]
\end{frame}

% 50 --------------------------------------------------------------------------
\begin{frame}{Две задачи с одним минимумом могут быть очень разными}
\small
Рассмотрим
\[
 f_1(x)=\frac12(x_1^2+x_2^2),
\]
\[
 f_2(x)=\frac12(x_1^2+100x_2^2).
\]
Обе функции:
\begin{itemize}
 \item положительно определены;
 \item имеют единственный минимум в $x^\star=0$;
 \item имеют очень простую явную формулу.
\end{itemize}

Но кривизна по разным направлениям во второй задаче отличается в $100$ раз. Нужна численная характеристика этой неоднородности.
\end{frame}

% 51 --------------------------------------------------------------------------
\begin{frame}{Число обусловленности}
\small
\begin{block}{Определение}
Для симметричной положительно определённой матрицы $A$ число
\[
 \boxed{
 \kappa(A)=\frac{\lammax(A)}{\lammin(A)}
 }
\]
называется спектральным числом обусловленности.
\end{block}

Так как $0<\lammin\le\lammax$,
\[
 \kappa(A)\ge1.
\]

\begin{alertblock}{Что именно измеряет $\kappa$?}
Не общую ``крутизну'' функции, а \textbf{неравномерность кривизны по разным направлениям}.
\end{alertblock}

Действительно, для любого $c>0$
\[
 \kappa(cA)=\frac{c\lammax(A)}{c\lammin(A)}=\kappa(A).
\]
\end{frame}

% 52 --------------------------------------------------------------------------
\begin{frame}{Число обусловленности --- отношение экстремальных кривизн}
\small
\begin{block}{Утверждение}
Если $A=A^\top\succ0$, то
\[
 \max_{\norm{v}=1}v^\top Av=\lammax(A),
 \qquad
 \min_{\norm{v}=1}v^\top Av=\lammin(A).
\]
Следовательно,
\[
 \boxed{
 \kappa(A)=
 \frac{\displaystyle\max_{\norm{v}=1}v^\top Av}
      {\displaystyle\min_{\norm{v}=1}v^\top Av}.
 }
\]
\end{block}

\begin{block}{Доказательство}
Из спектральных границ при $\norm{v}=1$:
\[
 \lammin\le v^\top Av\le\lammax.
\]
Обе границы достигаются на собственных векторах $v_{\min}$ и $v_{\max}$.
\end{block}
\end{frame}

% 53 --------------------------------------------------------------------------
\begin{frame}{Число обусловленности точно задаёт вытянутость линии уровня}
\footnotesize
\begin{columns}[T,onlytextwidth]
\column{0.37\textwidth}
В собственных координатах
\[
 \frac12\sum_i\lambda_i z_i^2=c.
\]
Полуось вдоль $v_i$ равна
\[
 a_i=\sqrt{\frac{2c}{\lambda_i}}.
\]
Поэтому
\[
 \frac{a_{\max}}{a_{\min}}
 =\sqrt{\frac{\lammax}{\lammin}}
 =\boxed{\sqrt{\kappa(A)}}.
\]

\begin{alertblock}{Точная геометрическая связь}
Если $\kappa=100$, то линия уровня вытянута по главным осям в $\sqrt{100}=10$ раз.
\end{alertblock}

\column{0.59\textwidth}
\centering
\includegraphics[width=0.98\linewidth]{condition_number_contours.pdf}
\end{columns}
\end{frame}

% 54 --------------------------------------------------------------------------
\begin{frame}{Большая кривизна сама по себе не означает плохую обусловленность}
\small
Сравним две матрицы:
\[
 A_1=I,
 \qquad
 A_2=100I.
\]
У второй функции кривизна во всех направлениях в $100$ раз больше, но
\[
 \kappa(A_1)=1,
 \qquad
 \kappa(A_2)=\frac{100}{100}=1.
\]

\begin{alertblock}{Ключевой смысл}
Хорошая обусловленность означает, что все главные направления имеют \textbf{сопоставимую} кривизну. Абсолютный масштаб собственных значений --- другой вопрос.
\end{alertblock}

\begin{block}{Самопроверка}
Что хуже обусловлено: $A=0.01I$ или $A=\operatorname{diag}(1,100)$?
\end{block}
\end{frame}

% 55 --------------------------------------------------------------------------
\begin{frame}{Почему это называется ``обусловленностью''?}
\footnotesize
Число $\kappa$ связано не только с формой линий уровня, но и с \textbf{чувствительностью решения} к небольшим изменениям задачи.

Пусть
\[
 x^\star=A^{-1}b,
\]
а после возмущения $b\mapsto b+\delta b$ решение изменилось на
\[
 x^\star\mapsto x^\star+\delta x.
\]
Тогда
\[
 A\,\delta x=\delta b.
\]

\begin{block}{Оценка чувствительности}
Для $b\ne0$
\[
 \boxed{
 \frac{\norm{\delta x}}{\norm{x^\star}}
 \le
 \kappa(A)
 \frac{\norm{\delta b}}{\norm{b}}.
 }
\]
\end{block}

То есть большое $\kappa$ допускает сильное усиление относительной ошибки входных данных.
\end{frame}

% 56 --------------------------------------------------------------------------
\begin{frame}{Откуда берётся оценка чувствительности?}
\footnotesize
Из $A\delta x=\delta b$ получаем
\[
 \delta x=A^{-1}\delta b.
\]
Собственные значения $A^{-1}$ равны $1/\lambda_i(A)$, поэтому
\[
 \norm{\delta x}
 \le
 \frac{1}{\lammin(A)}\norm{\delta b}.
\]
С другой стороны,
\[
 \norm{b}=\norm{Ax^\star}
 \le
 \lammax(A)\norm{x^\star},
\]
то есть
\[
 \frac1{\norm{x^\star}}
 \le
 \frac{\lammax(A)}{\norm{b}}.
\]
Перемножая оценки,
\[
 \frac{\norm{\delta x}}{\norm{x^\star}}
 \le
 \frac{\lammax}{\lammin}
 \frac{\norm{\delta b}}{\norm{b}}
 =
 \kappa(A)\frac{\norm{\delta b}}{\norm{b}}.
\]
\end{frame}

% 57 --------------------------------------------------------------------------
\begin{frame}{Одинаковое возмущение может по-разному сдвинуть минимум}
\small
Возьмём в обоих случаях
\[
 b=(0,1)^\top,
 \qquad
 \delta b=(0.1,0)^\top.
\]
Тогда относительное изменение $b$ равно $10\%$.

\begin{columns}[T,onlytextwidth]
\column{0.40\textwidth}
\begin{itemize}
 \item при $A=I$ имеем $\kappa=1$ и относительный сдвиг решения $10\%$;
 \item при $A=\operatorname{diag}(0.01,1)$ имеем $\kappa=100$ и относительный сдвиг решения $1000\%$.
\end{itemize}

Это пример, где оценка предыдущего слайда фактически достигается.

\column{0.56\textwidth}
\centering
\includegraphics[width=0.96\linewidth]{conditioning_sensitivity.pdf}
\end{columns}
\end{frame}

% 58 --------------------------------------------------------------------------
\begin{frame}{Почему обусловленность станет важна уже на следующей лекции}
\small
\begin{columns}[T,onlytextwidth]
\column{0.38\textwidth}
Градиент указывает локальное направление наискорейшего убывания:
\[
 -\grad f(x).
\]

Но при большом $\kappa$ поверхность имеет очень разные кривизны по главным направлениям. Градиент оказывается сильно ориентирован поперёк узкой долины.

\begin{alertblock}{Вопрос}
Как число обусловленности повлияет на допустимый шаг и скорость градиентного спуска?
\end{alertblock}

\column{0.58\textwidth}
\centering
\includegraphics[width=0.98\linewidth]{gradient_step_preview.pdf}
\end{columns}
\end{frame}

% 59 --------------------------------------------------------------------------
\begin{frame}{Карта всей лекции}
\small
\[
 \boxed{f(x)=\frac12x^\top Ax-b^\top x+c}
\]
\[
 \Downarrow
\]
\[
 \grad f(x)=Ax-b,
 \qquad
 \Hess f(x)=A
\]
\[
 \Downarrow
\]
\[
 A=Q\Lambda Q^\top
\]
\[
 \Downarrow
\]
\[
 \lambda_i
 \Longleftrightarrow
 \text{кривизна по главным направлениям}
\]
\[
 \Downarrow
\]
\[
 A\succ0
 \Longrightarrow
 x^\star=A^{-1}b,
 \qquad
 \kappa(A)=\frac{\lammax}{\lammin}.
\]
\end{frame}

% 60 --------------------------------------------------------------------------
\begin{frame}{Следующая лекция: почему работает градиентный спуск?}
\small
Теперь для квадратичной функции мы точно знаем
\[
 \grad f(x)=Ax-b.
\]
Градиентный шаг имеет вид
\[
 \boxed{
 x_{k+1}=x_k-\alpha(Ax_k-b).
 }
\]
А точное решение удовлетворяет
\[
 Ax^\star=b.
\]

\begin{alertblock}{Главный вопрос следующей лекции}
При каких значениях шага $\alpha$ последовательность $x_k$ действительно приходит к $x^\star$? И как скорость зависит от $\lammin$, $\lammax$ и $\kappa(A)$?
\end{alertblock}

Именно здесь обусловленность начнёт непосредственно управлять поведением алгоритма.
\end{frame}

\end{document}